
\documentclass[letter, 11pt]{article}

    %%%%%%%%%%%%%%%%%%%%%%%%%%%%%%
    \oddsidemargin  0.0in
    \evensidemargin 0.0in
    \textwidth      6.8in
    \headheight     28pt
    \topmargin      0.0in
    \textheight=9.0in
    \parindent=0in
    %%%%%%%%%%%%%%%%%%%%%%%%%%%%%%

\renewcommand{\familydefault}{\sfdefault} % Police 

\usepackage{amsthm,amsfonts,amssymb,amsmath, xcolor, enumitem}
\usepackage{graphicx}   %To include graphics
\usepackage[small,compact]{titlesec}

\usepackage{textcomp}
\newcommand{\mytilde}{$\sim$}
\usepackage{fancyvrb}
\fvset{framesep=1mm,fontfamily=courier,framerule=.3mm,numbersep=1mm,commandchars=\\\{\}}

\usepackage[T1]{fontenc}\usepackage{lmodern}\usepackage[utf8]{inputenc}

\usepackage[french]{babel}

\usepackage{fancyhdr}
\pagestyle{fancy}

\usepackage{version} % Pour faire un texte troué pour les étudiants
%\includeversion{etudiant}\excludeversion{professeur} %version etudiant
\includeversion{professeur}\excludeversion{etudiants} %version prof
\newcommand{\etudiant}[1]{\begin{etudiants}#1 \end{etudiants}}
\newcommand{\prof}[1]{\begin{professeur}\color{blue} {\bf Solution :} #1 \end{professeur}}


\fancyhead[L]{STT-2100 \\ Régression}
\fancyhead[R]{A.-S. Charest, E. Reny-Nolin, A.-A. Vallée \\ Université Laval}


\begin{document}

\vspace*{0.01 cm}
\begin{center}
{\Large Ajuster une RLS avec R}
\end{center}

Les mots en rouge indiquent les noms choisis par l'utilisateur.\\

Pour effectuer une analyse de régression linéaire simple avec R, on peut utiliser la fonction {\tt lm( )} (pour {\it linear model}), dont la syntaxe la plus simple prend la forme suivante :


\begin{Verbatim}[commandchars=\\\{\}]
{\textcolor{red}{mod}} <- lm({\textcolor{red}{varY}} \mytilde {\textcolor{red}{varX}}, data={\textcolor{red}{nomDuJeuDeDonnees}} )
\end{Verbatim}
où {\textcolor{red}{\texttt varY}} est la variable dépendante et {\textcolor{red}{\tt varX}} est la variable indépendante. 
Cette fonction retourne un objet de type liste et de classe lm, contenant plusieurs composantes dont les suivantes :
\begin{itemize}
\item \texttt{{\textcolor{red}{mod}}\$coefficients} : un vecteur contenant les coefficients ;
\item \texttt{{\textcolor{red}{mod}}\$residuals} : un vecteur contenant les résidus ;
\item \texttt{{\textcolor{red}{mod}}\$fitted.values} : un vecteur contenant les valeurs prédites.
\end{itemize}


La fonction \texttt{ls( )} appliquée à un objet de classe lm permet de connaître les noms de toutes ses composantes.
Pour accéder à une des composantes, on utilise {\tt \$} suivi du nom de la composante. Par exemple, {\tt {\textcolor{red}{mod}}\$coefficients} constitue le vecteur des coefficients $\hat\beta_0$ et $\hat\beta_1$ de la régression de {\textcolor{red}{\tt varY}} sur {\textcolor{red}{\tt varX}}.

\vspace{0.1in}
Pour tracer un nuage de points, on utilise la commande {\tt plot(x={\textcolor{red}{varX}},y={\textcolor{red}{varY}})}. Les paramètres graphiques peuvent être personnalisés, voir l'aide sur la commande {\tt par}.

\vspace{0.1in}
Pour ajouter une droite sur un graphique, utiliser la fonction \texttt{abline}, dont les arguments {\tt a} et  {\tt b} sont l’ordonnée à l’origine et la pente de la droite. Par exemple {\tt abline(a=0,b=1)} trace la droite $y=x$. On peut aussi ajouter la droite de régressionà l'aide de {\tt abline(reg={\textcolor{red}{mod}})}.


\end{document}

